%---------------------------------------------------------------------------%
%->> 项目自定义排版
%->  页面尺寸、章标题格式、图/表/公式标题样式、页眉页脚、参考文献与
%->  超链接均已由 Style/ucasthesis.cls 与 Style/artratex.sty 提供，
%->  本文件只补充本项目所需的宏与少量版式参数。
%---------------------------------------------------------------------------%
\usepackage{tcolorbox}
\usetikzlibrary{arrows.meta,positioning,fit,calc,backgrounds}
\usepackage{needspace}
%-\
%-> 图表与公式按章编号，与模板 Mainmatter 的约定一致
%-\
\renewcommand{\thefigure}{\thechapter-\arabic{figure}}
\renewcommand{\thetable}{\thechapter-\arabic{table}}
\renewcommand{\theequation}{\thechapter-\arabic{equation}}
%-\
%-> 正文段落与浮动体间距
%-\
\setlength{\parindent}{2em}
\setlength{\emergencystretch}{2em}
\setlength{\textfloatsep}{10pt plus 2pt minus 2pt}
\setlength{\floatsep}{9pt plus 2pt minus 2pt}
\setlength{\intextsep}{9pt plus 2pt minus 2pt}
%-\
%-> 表格列型
%-\
\newcolumntype{Y}{>{\raggedright\arraybackslash}X}
\newcolumntype{P}[1]{>{\raggedright\arraybackslash}p{#1}}
\renewcommand{\arraystretch}{1.22}
\setlength{\tabcolsep}{5pt}
%-\
%-> 颜色
%-\
\definecolor{ink}{HTML}{24364B}
\definecolor{muted}{HTML}{64748B}
\definecolor{codebg}{HTML}{FAF3EC}
\definecolor{codefg}{HTML}{8B2500}
\definecolor{asmbg}{HTML}{E3ECF7}
\definecolor{asmfg}{HTML}{0B3866}
\definecolor{softblue}{HTML}{EDF3F8}
\definecolor{softgreen}{HTML}{EDF5EE}
\definecolor{softamber}{HTML}{FBF3E4}
%-\
%-> 行内标记：\code 暖米底（标识符/函数名）、\asm 冷蓝底（汇编片段）、
%->          \codeplain 无底色（长片段，可断行）
%-\
\newcommand{\code}[1]{\tcbox[on line,colback=codebg,colframe=codebg,boxrule=0pt,
  arc=1.2pt,boxsep=1pt,left=1.5pt,right=1.5pt,top=0.5pt,bottom=0.5pt,
  fontupper=\ttfamily\small\color{codefg}]{\detokenize{#1}}}
\newcommand{\asm}[1]{\tcbox[on line,colback=asmbg,colframe=asmfg!35,boxrule=0.3pt,
  arc=1.2pt,boxsep=1pt,left=1.5pt,right=1.5pt,top=0.5pt,bottom=0.5pt,
  fontupper=\ttfamily\small\color{asmfg}]{#1}}
\newcommand{\codeplain}[1]{\texttt{\detokenize{#1}}}
%-\
%-> 代码块：中文标签、细小字号、语法高亮
%-\
%  artratex 在检测到 XITS 时把等宽族也指向 XITS，而 XITS 缺 U+2192/U+2191
%  等符号字形，代码块里的箭头会报 Missing character。这里把等宽族改回
%  DejaVu Sans Mono（含箭头字形），中文由 CJKmonofont 兜底。
\setmonofont{DejaVu Sans Mono}[Scale=0.82]
\renewcommand{\lstlistingname}{代码块}
\lstset{%
  numbers=none, title={}, stepnumber=1, numberstyle=\tiny\color{muted},
  basicstyle=\small\ttfamily, breaklines=true, columns=flexible,
  keepspaces=true, frame=single, rulecolor=\color{black!25},
  showstringspaces=false, backgroundcolor=\color{white}, tabsize=4,
  aboveskip=7pt, belowskip=7pt, captionpos=b,
  language=C,
  keywordstyle=\color{blue!60!black}\bfseries,
  commentstyle=\color{muted}\itshape,
  stringstyle=\color{red!55!black},
  emph={translator_io_start,gen_helper_csrrw,gen_set_gpr,do_csr_post,dest_gpr,
    tcg_constant_i32,HELPER,riscv_csrrw,riscv_raise_exception,decode_save_opc,
    gen_update_pc,exit_tb,DEF_HELPER_4},
  emphstyle=\color{teal!60!black}\bfseries,
  morekeywords={uint64_t,uint32_t,uint8_t,int8_t,target_ulong,TCGv,TCGv_i32,
    TCGv_i64,DisasContext,DisasContextBase,CPURISCVState,RISCVException,
    CPUArchState,privilege_mode_t,bool,true,false}}
%-\
%-> Python 伪代码样式：listings 的 Python 只有 [2]/[3] 两个方言，必须显式选 3
%-\
\lstdefinestyle{pseudopy}{
  language={[3]Python},
  keywordstyle=\color{blue!60!black}\bfseries,
  commentstyle=\color{muted}\itshape,
  stringstyle=\color{red!55!black},
  emphstyle=\color{teal!60!black}\bfseries,
  emph={cpu_exec_loop,tb_gen_code,cpu_handle_interrupt,tcg_handle_interrupt,
    address_space_dispatch,virtio_blk_handle_vq,guest_load,device_raise_irq,
    kvm_arch_vcpu_ioctl_run,check_vcpu_requests,kvm_vgic_flush_hwstate,
    kvm_arm_vcpu_enter_exit,handle_exit,smp_store_mb,kvm_vgic_sync_hwstate},
  morekeywords=[2]{None,True,False}
}
%-\
%-> 机制示意图样式
%-\
\tikzset{box/.style={draw=ink!70,rounded corners=2pt,fill=softblue,align=center,
    font=\small,minimum height=10mm,inner sep=5pt},
  greenbox/.style={box,fill=softgreen},amberbox/.style={box,fill=softamber},
  wire/.style={-{Stealth[length=2mm]},draw=ink,thick},
  dwire/.style={wire,dashed},
  note/.style={font=\footnotesize,align=center,text=muted}}
%-\
%-> 封面：课程调研报告版
%->  模板的学位论文封面带有指导教师、学位、培养单位等信息，
%->  本报告是课程作业而非学位申请，故只保留作者、学号与课程。
%-\
\newcommand{\makecoverpage}{%
  \cleardoublepage
  \thispagestyle{empty}%
  \begin{center}%
    \linespread{1.6}\zihao{4}\bfseries

    \vspace*{\stretch{4}}

    {\includegraphics[scale=0.095]{ucas_logo}}

    \vspace*{\stretch{2}}

    {\zihao{1}\bfseries\sffamily 课程调研报告}

    \vspace*{\stretch{3}}

    {\zihao{-2}\bfseries\sffamily \ulenhance[1.5pt]{\ulextend{虚拟化技术}}}

    \vspace*{1.6ex}

    \vspace*{\stretch{3}}

    \def\tabcolsep{1pt}\def\arraystretch{1.35}%
    \begin{tabular}{lc}%
      作者姓名 & \ulenhance[1.2pt]{\ulextend{王\quad 靳}}\\
      学\qquad 号 & \ulenhance[1.2pt]{\ulextend{2026E8013282153}}\\
      课\qquad 程 & \ulenhance[1.2pt]{\ulextend{高级操作系统教程}}\\
    \end{tabular}

    \vspace*{\stretch{4}}

    {2026~年~11~月}

    \vspace*{\stretch{4}}
  \end{center}%
  \clearpage
  \cleardoublepage[empty]%
}
%-\
%-> 长 URL 可断行（hyperref 由 artratex 在 \AtEndPreamble 中加载）
%-\
\AtEndPreamble{\usepackage{xurl}}
%---------------------------------------------------------------------------%
